\input{../article_base.tex}
\title{פתרון מטלה 11 --- משוואות דיפרנציאליות, 80320}
% chktex-file 9
% chktex-file 17

\DeclareMathOperator\Div{div}

\usepackage{pgfplots, tikz, pgffor}
\usetikzlibrary{math}
\pgfplotsset{compat=1.18}

\begin{document}
\maketitle
\maketitleprint[blue]{}

\question{}
נגדיר את $\Phi : \RR^2 \setminus \{ 0 \} \to \RR$ על־ידי,
\[
	\Phi(x)
	= \frac{1}{2 \pi} \ln \lVert x \rVert
\]
נבחין כי מצאנו שזהו פתרון למשוואת לפלס הדו־מימדית.
נראה שאם $f \in C_c^2(\RR^2)$ הפונקציה,
\[
	u(x)
	= \frac{1}{2 \pi} \int_{\RR^2} \Phi(y) f(x - y)\ dy
	= \frac{1}{2 \pi} (\Phi * f)(x)
\]
היא פתרון למשוואה $\Delta u = f$.
\begin{proof}
	נראה תחילה שהיא מוגדרת לכל $x$.
	מחסימות וקומפקטיות $f$ מספיק לבדוק את התכנסות האינטגרל של $\Phi$ בכדורים פתוחים סביב הראשית,
	\[
		\int_{B(0, R)} \frac{1}{2 \pi} \ln \lVert x \rVert\ dx
		= \int_0^{2 \pi} \int_0^R \frac{1}{2 \pi} r \ln(r)\ dr\ d \theta
		= \frac{1}{4} \left. x^2 2 \ln x - x^2 \right\rvert_0^R
		< \infty
	\]
	וקיבלנו כי $u$ אכן מוגדרת, ונשאר להראות שהיא מקיימת $\Delta u = f$.
	נבחין כי מהמטלה הקודמת נובע,
	\[
		\Delta u
		= \frac{1}{2 \pi} \int_{\RR^2} \Phi(y) \Delta f(y - x)\ dy
	\]
	בפרט $\Delta u$ מוגדר ונוכל להוכיח את הטענה.
	נבחין שבשונה מהמקרה $d > 2$ האינטגרל מתכנס רק בכדורים סביב הראשית ולכן נגדיר את $\Omega_r = B(0, r)$ ונקבל,
	\[
		\Delta u
		= \frac{1}{2 \pi} \int_{\RR^2} \Phi(y) \Delta f(y - x)\ dy
		= \frac{1}{2 \pi} \lim_{r \to \infty} \int_{\Omega_r} \Phi(y) \Delta f(y - x)\ dy
	\]
	בתחום זה מנוסחת גרין הראשונה,
	\[
		\int_{\Omega_r} \Phi(y) \Delta f(y - x)\ dy
		= \int_{\partial B(0, r)} \Phi(y) \partial_{\nu} f(x - y)\ dy - \int_{B(0, r)} \langle \nabla \Phi(y), \nabla f(x, y) \rangle\ dy
	\]
	אנו יודעים ש־$f |_{B^C(0, r)} \equiv 0$ עבור $r$ גדול מספיק, לכן בפרט נגזרתה הנורמלית מתאפסת בסביבה זו ונקבל,
	\[
		\int_{\Omega_r} \Phi(y) \Delta f(y - x)\ dy
		= - \int_{B(0, r)} \langle \nabla \Phi(y), \nabla f(x, y) \rangle\ dy
	\]
	וממהלך שקול לתרגול על־ידי שימוש בנוסחת גרין ובהרמוניות $\Phi$ נקבל,
	\[
		\int_{\Omega_r} \Phi(y) \Delta f(y - x)\ dy
		= \int_{\partial B(0, r)} \partial_{\nu} \Phi(y) f(x - y)\ dy
	\]
	אבל,
	\[
		\partial_{\nu} \Phi(y)
		= \frac{1}{2 \pi} \frac{y}{{\lVert y \rVert}^2}
	\]
	וכן כפונקציה רדיאלית $\partial_{\nu} \Phi$ קבועה ונוכל לחשבה, מערך הנגזרת נוכל להסיק,
	\[
		\int_{\partial B(0, r)} \partial_{\nu} \Phi(y)\ dy
		= \frac{1}{r} \frac{1}{2 \pi} \vol(\partial B(0, r))
		= 1
	\]
	וקיבלנו את,
	\[
		\int_{\partial B(0, r)} f(x - y)\ dy
	\]
	ומרציפות הגבול הזה הוא בדיוק $f(x)$ וסיימנו.
\end{proof}

\question{}
\subquestion{}
תהי $\Omega \subseteq \RR^d$ פתוחה וחסומה בעלת שפה חלקה.
נניח ש־$u_1, u_2 \in C^2(\overline{\Omega})$ כך שלכל $x \in \Omega$ מתקיים,
\[
	\Delta u_1(x)
	= \Delta u_2(x)
\]
ונניח בנוסף ש־$u_1 |_{\partial \Omega} \equiv u_2 |_{\partial \Omega} \equiv 0$.
נוכיח ש־$u_1 = u_2$.
\begin{proof}
	נסמן $u = u_1 - u_2$.
	נזכור כי לכל $f \in C^1(\Omega)$ מתקיים,
	\[
		\Div(f \cdot \nabla f)
		= \sum_{n = 1}^d \partial_n (f \cdot \partial_n f)
		= \sum_{n = 1}^d {(\partial_n f)}^2 + f \cdot \partial_{n^2} f
		= \lVert \nabla f \rVert^2 + f \cdot \nabla f
	\]
	נזכור כי מתקיים $u |_{\partial \Omega} = (u_1 - u_2) |_{\partial \Omega} = 0$ ולכן ממשפט הדיברגנץ מתקיים,
	\[
		\int_{\Omega} \lVert \nabla u \rVert^2 + u \cdot \Delta u\ d\vol_d
		= \int_{\Omega} \Div(u \cdot \nabla u)\ d\vol_d
		= \int_{\partial \Omega} \langle u \cdot \nabla u, \nu \rangle\ d\vol_{d - 1}
		= 0
	\]
	אבל נזכור כי גם $\Delta u_1 = \Delta u_2 \implies \Delta u = 0$ וקיבלנו שבדיוק,
	\[
		\int_{\Omega} \lVert \nabla u \rVert^2\
		= 0
	\]
	ובהתאם נסיק ש־$u = 0$ בלבד (זאת שכן זהו פונקציונל לינארי חיובי) ובפרט $u_1 = u_2$.
\end{proof}

\subquestion{}
נבחין כי הסעיף הקודם נכון גם תחת ההנחה היותר כללית ש־$\partial_{\nu} u_1 = \partial_{\nu} u_2$ על $\partial \Omega$,
זאת שכן זאת הייתה הדרישה היחידה שהשתמשנו בה במהלך ההוכחה.

\subquestion{}
נניח ש־$u_1, u_2 \in C(\RR^d)$ כך שלכל $x \in \RR^d$ מתקיים,
\[
	\Delta u_1(x)
	= \Delta u_2(x)
\]
וכן נניח ש־$\lim_{x \to \infty} u_1(x) - u_2(x) = 0$. \\
נראה ש־$u_1 = u_2$.
\begin{proof}
	נסמן שוב $u = u_1 - u_2$ ולכן מהנתונים נובע $\Delta u \equiv 0$ וכן $\lim_{x \to \infty} u(x) = 0$.
	קיבלנו ש־$u$ היא הרמונית ומקיימת את תכונת הערך הממוצע על ספירות, אבל לכל $\varepsilon > 0$ קיים $r > 0$ כך שאם $x \in B(\infty, r)$ אז,
	\[
		\lVert u(x) \rVert < \varepsilon
	\]
	ישירות מהגדרת הגבול הנתון,
	לכל $x \in \RR^d$ נוכל אם כך לבחור ספירה עם רדיוס $r$ גדול מספיק ונקבל מעיקרון המקסימום והמינימום ומתכונת הערך הממוצע על אינטגרל שמתקיים,
	\[
		- \vol_{d - 1}(\partial B(x, r)) \varepsilon
		\le u(x)
		= \int_{\partial B(x, r)} u(y)\ dy
		\le \vol_{d - 1}(\partial B(x, r)) \varepsilon
	\]
	ומשרירותיות $\varepsilon > 0$ נקבל שבפרט $u(x) = 0$, כלומר $u_1 = u_2$ על $\RR^d$.
\end{proof}

\question{}
תהי $u : \RR^d \to \RR$ הרמונית.
נראה שאם $u$ חסומה אז $u$ קבועה.
\begin{proof}
	ראינו בכיתה שקיים קבוע $0 < C = C(d, 1)$ כך שמתקיים,
	\[
		\forall x \in \RR^d,\ 
		|\partial_i u| \le \frac{C}{r^{d + 1}} \int_{B(x, r)} |u(x)|\ dx
	\]
	נסמן $0 < D$ חסם של $u$, כלומר $|u| \le D$ ולכן,
	\[
		|\partial_i u|
		\le \frac{C}{r^{d + 1}} \int_{B(x, r)} |u(x)|\ dx
		\le \frac{C}{r^{d + 1}} \vol(B(0, r)) D
	\]
	אבל $\vol_d(B(0, r)) = o(r^d)$ ולכן נקבל שעבור $r \to \infty$ מתקבל $|\partial_i u| \le 0$, כלומר $D u \equiv 0$ ולכן $u$ קבועה.
\end{proof}

\question{}
יהי $d \ge 3$ ונגדיר $K : B(0, 1) \times \partial B(0, 1)$ על־ידי,
\[
	K(x, y)
	= \frac{1}{\vol_{d - 1}(\partial B_1)}  \frac{1 - \lVert x \rVert^2}{\lVert x - y \rVert^2}
\]
ותהי $g : \partial B(0, 1) \to \RR$ רציפה.

\subquestion{}
נראה שהפונקציה,
\[
	u(x)
	= \int_{\partial B(0, 1)} K(x, y) g(y)\ ds(y)
\]
היא הרמונית.
\begin{proof}
	נראה ש־$u$ מקיימת את תכונת הערך הממוצע לספירות.
	יהיו $x_0 \in B(0, 1), r > 0$ כך ש־$B(x_0, r) \subseteq B(0, 1)$ ונבדוק,
	\begin{multline*}
		\int_{B(x_0, r)} u(x)\ dx
		= \int_{B(x_0, r)} \int_{\partial B(0, 1)} K(x, y) g(y)\ ds(y)\ dx \\
		= \int_{\partial B(0, 1)} g(y) \int_{B(x_0, r)} K(x, y)\ dx\ ds(y)
		= \int_{\partial B(0, 1)} g(y) K(x_0, y)\ ds(y)
		= u(x_0)
	\end{multline*}
	כאשר השתמשנו בפוביני ובעובדה שלכל $y$ הפונקציה $K(x, y)$ היא הרמונית ביחס ל־$x$.
	בהתאם $u$ הרמונית.
\end{proof}

\subquestion{}
נראה שלכל $y_0 \in \partial B(0, 1)$ מתקיים,
\[
	\lim_{x \to x_0} u(x)
	= g(y_0)
\]
\begin{proof}
	יהי $\varepsilon > 0$.
	יהיו $x \in B(0, 1), y_0 \in \partial B(0, 1)$ כך שמרחקן קטן מספיק.
	
	קיימת קשת $A \subseteq \partial B(0, 1)$ כך שמתקיים,
	\[
		\left\lvert \int_{\partial B(0, 1) \setminus A} K(x, y) g(y)\ ds(y) \right\rvert
		< \varepsilon
	\]
	זאת ישירות מתכונת הערך הממוצע והחסרת קשת מספיק קטנה.
	נוכל להקטין מספיק עוד את $A$ כדי לקבל מרציפות $g$ את הטענה,
	\[
		\sup_{y \in A} |g(y) - g(y_0)|
		< \varepsilon
	\]
	מתקיים,
	\[
		|u(x) - g(y_0)|
		= \left\lvert u(x) - \int_{\partial B(0, 1)} K(x, y) g(y_0)\ ds(y) \right\rvert
		=
		\left\lvert \int_{\partial B(0, 1)} K(x, y) (\overbrace{g(y) - g(y_0)}^{\le \varepsilon})\ ds(y) \right\rvert
		\le \varepsilon
	\]
	כאשר השתמשנו בעובדה ש־$\int K = 1$.
	בהתאם נקבל שבדיוק $\lim_{x \to y_0} u(x) = g(y_0)$.
\end{proof}

\end{document}
